% 執筆者一覧：本書に関わった人を並べる
% 行を追加するときは、tabular の中に「名前 & 担当 \\」の形で 1 行ずつ書き足す。
\clearpage\phantomsection\addcontentsline{toc}{part}{執筆者一覧}
\chapter*{執筆者一覧}
\markboth{執筆者一覧}{}

本書は、ロボットチーム eR@sers のメンバーによって執筆されました。

\begin{tablebox}
  \renewcommand{\arraystretch}{1.3}
  \begin{tabular}{lp{0.6\linewidth}}
    \toprule
    名前 & 担当 \\
    \midrule
    坂巻 新 & 全体の構成、執筆 \\
    % 名前 & 担当 \\
    \bottomrule
  \end{tabular}
\end{tablebox}

\vspace{2\baselineskip}

\begin{description}
  \item[本書の原稿] \url{https://github.com/trcp/erasers_book}
  \item[サンプルコード] \codeurl
\end{description}
